\section{Formulación Matemática del Modelo Determinista de Distribución}

\subsection{Conjuntos y Índices}

\begin{itemize}
    \item $K$: conjunto de camiones disponibles, $k \in K$.
    \item $C$: conjunto de compartimentos de cada camión, $c \in C = \{C_0, C_1\}$.
    \item $P$: conjunto de productos, $p \in P = \{R, D\}$, donde $R$ representa Regular y $D$ representa Diésel.
    \item $N$: conjunto de estaciones de servicio, $j \in N$.
    \item $V$: conjunto de todos los nodos, $V = N \cup \{0\}$, donde $0$ representa el depósito central.
\end{itemize}

\subsection{Parámetros}

\begin{itemize}
    \item $d_{ij}$: distancia entre el nodo $i$ y el nodo $j$ [km].
    \item $D_{jp}$: demanda de la estación $j$ para el producto $p$ [L].
    \item $Q_{kc}$: capacidad del compartimento $c$ del camión $k$ [L].
    \item $F_k$: costo fijo por utilizar el camión $k$.
    \item $c_d$: costo por kilómetro recorrido.
    \item $c_s$: penalización por litro de combustible no entregado.
    \item $a_j$: inicio de la ventana de tiempo de la estación $j$.
    \item $b_j$: término de la ventana de tiempo de la estación $j$.
    \item $T_j$: tiempo de servicio en la estación $j$.
    \item $v_k$: velocidad promedio del camión $k$.
    \item $h_k$: hora de salida programada del camión $k$ desde el depósito.
    \item $\Delta$: máxima diferencia permitida entre los \textit{fill ratios}, con $\Delta = 0.30$.
    \item $M$: constante suficientemente grande para modelar restricciones lógicas.
\end{itemize}

El costo fijo $F_k$ se incluye en la función objetivo debido a que se entrega directamente este valor en la tabla de características de la flota. Asimismo, la hora de salida $h_k$ se considera como una hora programada para el ruteo, mientras que en la etapa de scheduling se verifica si los camiones pueden estar cargados y listos antes de dichas horas de salida.

\subsection{Variables de Decisión}

\begin{itemize}
    \item $x_{ijk} \in \{0, 1\}$: vale 1 si el camión $k$ viaja directamente desde el nodo $i$ al nodo $j$, y 0 en caso contrario.
    \item $y_k \in \{0, 1\}$: vale 1 si el camión $k$ es utilizado, y 0 en caso contrario.
    \item $w_{jk} \in \{0, 1\}$: vale 1 si la estación $j$ es visitada por el camión $k$, y 0 en caso contrario.
    \item $z_{kcp} \in \{0, 1\}$: vale 1 si el compartimento $c$ del camión $k$ transporta el producto $p$, y 0 en caso contrario.
    \item $q_{kcpj} \geq 0$: cantidad de producto $p$ entregada por el compartimento $c$ del camión $k$ a la estación $j$.
    \item $s_{jp} \geq 0$: demanda no satisfecha del producto $p$ en la estación $j$.
    \item $L_{kci} \geq 0$: carga restante en el compartimento $c$ del camión $k$ después de visitar el nodo $i$.
    \item $t_{jk} \geq 0$: hora de llegada del camión $k$ a la estación $j$.
    \item $u_{jk} \geq 0$: variable auxiliar para la eliminación de subtours.
\end{itemize}

\subsection{Función Objetivo}

El objetivo es minimizar el costo total de operación, compuesto por el costo de distancia recorrida, el costo fijo de los camiones utilizados y la penalización por demanda no satisfecha.

\begin{equation}
    \min Z =
    c_d \sum_{k \in K} \sum_{i \in V} \sum_{\substack{j \in V \\ j \neq i}} d_{ij} x_{ijk}
    +
    \sum_{k \in K} F_k y_k
    +
    c_s \sum_{j \in N} \sum_{p \in P} s_{jp}
\end{equation}

\subsection{Restricciones}

\subsubsection{Ruteo y Conservación de Flujo}

Estas restricciones controlan el movimiento de los camiones dentro de la red de distribución.

\paragraph{Salida y retorno al depósito.}

Si un camión $k$ es utilizado, debe salir desde el depósito y regresar a él al finalizar su ruta.

\begin{align}
    \sum_{j \in N} x_{0jk} &= y_k, && \forall k \in K \\
    \sum_{i \in N} x_{i0k} &= y_k, && \forall k \in K
\end{align}

\paragraph{Conservación de flujo en estaciones.}

Si un camión entra a una estación, también debe salir de ella.

\begin{align}
    \sum_{\substack{i \in V \\ i \neq j}} x_{ijk} &= w_{jk}, && \forall j \in N, \forall k \in K \\
    \sum_{\substack{h \in V \\ h \neq j}} x_{jhk} &= w_{jk}, && \forall j \in N, \forall k \in K
\end{align}

\paragraph{Unicidad de visita.}

Cada estación puede ser visitada como máximo por un camión.

\begin{equation}
    \sum_{k \in K} w_{jk} \leq 1,
    \quad \forall j \in N
\end{equation}

\subsubsection{Satisfacción de Demanda y Capacidad}

\paragraph{Demanda con shortage.}

La demanda de cada estación y producto debe ser satisfecha mediante entregas o registrada como demanda no satisfecha.

\begin{equation}
    \sum_{k \in K} \sum_{c \in C} q_{kcpj} + s_{jp} = D_{jp},
    \quad \forall j \in N, \forall p \in P
\end{equation}

\paragraph{Compatibilidad compartimento-producto.}

Cada compartimento puede transportar como máximo un tipo de combustible.

\begin{equation}
    \sum_{p \in P} z_{kcp} \leq 1,
    \quad \forall k \in K, \forall c \in C
\end{equation}

\paragraph{Asignación de entrega.}

Solo se puede entregar un producto desde un compartimento si dicho producto fue asignado al compartimento y si el camión visita la estación correspondiente.

\begin{align}
    q_{kcpj} &\leq M z_{kcp},
    && \forall k \in K, \forall c \in C, \forall p \in P, \forall j \in N \\
    \sum_{c \in C} \sum_{p \in P} q_{kcpj}
    &\leq M w_{jk},
    && \forall k \in K, \forall j \in N
\end{align}

\paragraph{Límite de capacidad.}

La carga total asignada a cada compartimento no puede superar su capacidad.

\begin{equation}
    \sum_{j \in N} \sum_{p \in P} q_{kcpj} \leq Q_{kc},
    \quad \forall k \in K, \forall c \in C
\end{equation}

\subsubsection{Estabilidad de Carga}

\paragraph{Carga inicial.}

La carga inicial de cada compartimento corresponde a la suma de las cantidades que dicho compartimento entregará durante la ruta.

\begin{equation}
    L_{kc0} =
    \sum_{j \in N} \sum_{p \in P} q_{kcpj},
    \quad \forall k \in K, \forall c \in C
\end{equation}

\paragraph{Actualización de carga.}

Si el camión $k$ viaja desde el nodo $i$ hacia la estación $j$, la carga restante después de visitar $j$ debe ser igual a la carga restante después de visitar $i$, menos la cantidad entregada en $j$.

\begin{align}
    L_{kcj}
    &\leq
    L_{kci}
    -
    \sum_{p \in P} q_{kcpj}
    +
    M(1-x_{ijk}),
    && \forall k \in K, \forall c \in C, \forall i \in V, \forall j \in N, i \neq j \\
    L_{kcj}
    &\geq
    L_{kci}
    -
    \sum_{p \in P} q_{kcpj}
    -
    M(1-x_{ijk}),
    && \forall k \in K, \forall c \in C, \forall i \in V, \forall j \in N, i \neq j
\end{align}

\paragraph{Estabilidad en ruta.}

La diferencia entre los \textit{fill ratios} de los compartimentos de un mismo camión no puede superar $\Delta = 0.30$ después de visitar cada estación.

\begin{align}
    \frac{L_{kcj}}{Q_{kc}}
    -
    \frac{L_{kc'j}}{Q_{kc'}}
    &\leq
    \Delta + M(1-w_{jk}),
    && \forall k \in K, \forall c,c' \in C, c \neq c', \forall j \in N \\
    \frac{L_{kc'j}}{Q_{kc'}}
    -
    \frac{L_{kcj}}{Q_{kc}}
    &\leq
    \Delta + M(1-w_{jk}),
    && \forall k \in K, \forall c,c' \in C, c \neq c', \forall j \in N
\end{align}

\paragraph{Estabilidad inicial en depósito.}

La estabilidad también debe cumplirse al momento de salir desde el depósito.

\begin{align}
    \frac{L_{kc0}}{Q_{kc}}
    -
    \frac{L_{kc'0}}{Q_{kc'}}
    &\leq
    \Delta,
    && \forall k \in K, \forall c,c' \in C, c \neq c' \\
    \frac{L_{kc'0}}{Q_{kc'}}
    -
    \frac{L_{kc0}}{Q_{kc}}
    &\leq
    \Delta,
    && \forall k \in K, \forall c,c' \in C, c \neq c'
\end{align}

\subsubsection{Ventanas de Tiempo y Subtours}

\paragraph{Ventanas de tiempo.}

Si una estación es visitada, la hora de llegada debe encontrarse dentro de su ventana de tiempo.

\begin{align}
    t_{jk} &\geq a_j - M(1-w_{jk}),
    && \forall j \in N, \forall k \in K \\
    t_{jk} &\leq b_j + M(1-w_{jk}),
    && \forall j \in N, \forall k \in K
\end{align}

\paragraph{Propagación temporal.}

La llegada a una estación debe considerar la llegada al nodo anterior, el tiempo de servicio y el tiempo de viaje.

\begin{align}
    t_{jk}
    &\geq
    t_{ik} + T_i + \frac{d_{ij}}{v_k}
    -
    M(1-x_{ijk}),
    && \forall k \in K, \forall i,j \in N, i \neq j \\
    t_{jk}
    &\geq
    h_k + \frac{d_{0j}}{v_k}
    -
    M(1-x_{0jk}),
    && \forall k \in K, \forall j \in N
\end{align}

\paragraph{Eliminación de subtours.}

Se incorporan restricciones tipo MTZ para evitar ciclos desconectados del depósito.

\begin{align}
    u_{ik} - u_{jk} + |N|x_{ijk}
    &\leq |N|-1,
    && \forall k \in K, \forall i,j \in N, i \neq j \\
    1 \leq u_{jk}
    &\leq |N|w_{jk},
    && \forall j \in N, \forall k \in K
\end{align}